% Options for packages loaded elsewhere
\PassOptionsToPackage{unicode}{hyperref}
\PassOptionsToPackage{hyphens}{url}
\PassOptionsToPackage{dvipsnames,svgnames,x11names}{xcolor}
\documentclass[
  10pt,
  fontset=fandol]{ctexart}
\usepackage{xcolor}
\usepackage[margin=16mm]{geometry}
\usepackage{amsmath,amssymb}
\setcounter{secnumdepth}{-\maxdimen} % remove section numbering
\usepackage{iftex}
\ifPDFTeX
  \usepackage[T1]{fontenc}
  \usepackage[utf8]{inputenc}
  \usepackage{textcomp} % provide euro and other symbols
\else % if luatex or xetex
  \usepackage{unicode-math} % this also loads fontspec
  \defaultfontfeatures{Scale=MatchLowercase}
  \defaultfontfeatures[\rmfamily]{Ligatures=TeX,Scale=1}
\fi
\usepackage{lmodern}
\ifPDFTeX\else
  % xetex/luatex font selection
\fi
% Use upquote if available, for straight quotes in verbatim environments
\IfFileExists{upquote.sty}{\usepackage{upquote}}{}
\IfFileExists{microtype.sty}{% use microtype if available
  \usepackage[]{microtype}
  \UseMicrotypeSet[protrusion]{basicmath} % disable protrusion for tt fonts
}{}
\makeatletter
\@ifundefined{KOMAClassName}{% if non-KOMA class
  \IfFileExists{parskip.sty}{%
    \usepackage{parskip}
  }{% else
    \setlength{\parindent}{0pt}
    \setlength{\parskip}{6pt plus 2pt minus 1pt}}
}{% if KOMA class
  \KOMAoptions{parskip=half}}
\makeatother
\setlength{\emergencystretch}{3em} % prevent overfull lines
\providecommand{\tightlist}{%
  \setlength{\itemsep}{0pt}\setlength{\parskip}{0pt}}
\usepackage{amsmath,amssymb,mathtools,needspace}
\xeCJKDeclareCharClass{CJK}{"2032,"0400->"04FF,"2160->"216B,"2460->"2473,"25A0,"25C7,"25CB,"25CF}
\IfFontExistsTF{Arial Unicode MS}{\xeCJKsetup{AutoFallBack=true}\setCJKfallbackfamilyfont{\CJKrmdefault}{Arial Unicode MS}}{}
\setcounter{secnumdepth}{0}
\setlength{\emergencystretch}{2em}
\allowdisplaybreaks
\usepackage{bookmark}
\IfFileExists{xurl.sty}{\usepackage{xurl}}{} % add URL line breaks if available
\urlstyle{same}
\hypersetup{
  pdftitle={小结},
  colorlinks=true,
  linkcolor={Maroon},
  filecolor={Maroon},
  citecolor={Blue},
  urlcolor={Blue},
  pdfcreator={LaTeX via pandoc}}

\title{小结}
\author{}
\date{}

\begin{document}
\maketitle

\subsection{小结}\label{ux5c0fux7ed3}

方程 \(f(x)=0\) 求根有多种方法，本章重点介绍了 Newton 法。Newton 法是一种行之有效的迭代法，在单根附近具有较高的收敛速度。应用 Newton 法的关键在于选取足够精确的初值。如果初值选取不当（偏离所求的根比较远），则 Newton 法可能发散。Newton 下山法（见 6.3.5 节）有时可用来克服这个局限性。

Newton 法的另一个局限性是要求计算导数值。如果函数 \(f\) 的形式复杂而不便于求导，则可用导数的估值或差商代替导数，而得出简化的 Newton 法（见 6.3.1 节）或近似的 Newton 法（见 6.4.1 节）。不过这样处理势必会影响收敛速度。

应当指出，对于代数方程 \(a_0x^n+a_1x^{n-1}+\cdots+a_n=0\)，其求根问题包含有许多深刻的理论和有效的算法，关于这方面的知识，请读者参看文献 {[}9{]} 的第 6 章《高次代数方程解法》。

\begin{center}\rule{0.5\linewidth}{0.5pt}\end{center}

\subsection{相关原页校注（agent 补充）}\label{ux76f8ux5173ux539fux9875ux6821ux6ce8agent-ux8865ux5145}

以下校注来自本节覆盖的原页，可能也涉及同页相邻小节；原文照录与校正说明分开保留。

\subsubsection{原书 PDF 页 175 的校注}\label{ux539fux4e66-pdf-ux9875-175-ux7684ux6821ux6ce8}

\href{../pages/pdf-175.md}{完整原页转录} · \href{../../source-images/pdf-175.jpeg}{原页图像}

校注记录：ch06-E006。

\begin{quote}
校注（agent 补充）：习题 2 的区间在原页确为 \((0,1)\)，转录保留。对 \(0<x<1\) 有 \(0<x\sin x<1\)，从而 \(1-x\sin x>0\)，故该区间不含题述方程的根。此为原题条件的疑误；未替教材指定修正后的区间，也未代做该题。
\end{quote}

\subsection{本节覆盖原页的校注记录}\label{ux672cux8282ux8986ux76d6ux539fux9875ux7684ux6821ux6ce8ux8bb0ux5f55}

下列记录按原页关联，含同页相邻小节的校注；计算时同时读取上文校注和完整原页。

\begin{itemize}
\tightlist
\item
  \texttt{ch06-E006}：原书 PDF 页 175
\end{itemize}

\end{document}
